\documentclass{article}
\usepackage{graphicx} % Required for inserting images

\title{Lab3 MIRPR}
\author{Emanuel Marcu}

\begin{document}

\maketitle

\section{Descrierea setului real de date si scurta EDA}

Setul de date folosit este compus din mai multe fisiere text, fiecare reprezentand o boala stomatologica (ex: abces parodontal, carie simpla, pericoronarita, pulpita etc.). 
Fiecare fisier contine aproximativ 30 de conversatii intre Student si Pacient, in limbaj informal.

\subsection*{Sinteza EDA}
\begin{itemize}
    \item conversatiile sunt separate prin ``CONVERSATIA X''
    \item replicile sunt sub forma ``Student:'' si ``Pacient:''
    \item limbajul pacientului este colocvial, natural si uneori incomplet
    \item pe baza conversatiilor am extras manual setul de simptome pentru fiecare boala
    \item simptomele extrase au devenit baza interna simptome $\rightarrow$ boala folosita la generarea prompturilor
\end{itemize}


\section{Metodologie experimentala}

\subsection*{Etapele principale}
\begin{itemize}
    \item citirea fisierelor si extragerea replicilor Student/Pacient
    \item conversia replicilor in format \textit{user/assistant}
    \item generarea unui mesaj de sistem ce include:
    \begin{itemize}
        \item boala reala (doar intern pentru model)
        \item lista simptomelor extrase manual pentru acea boala
    \end{itemize}
    \item formatarea conversatiilor cu \texttt{apply\_chat\_template}
    \item tokenizare la maxim 2048 tokeni
    \item folosirea tehnicii \textit{transfer learning} pentru specializarea modelelor pe datele noastre
    \item finetuning separat pentru LLaMA si RoLLaMA
\end{itemize}


\section{Impartirea setului de date}

\begin{itemize}
    \item 90\% -- train
    \item 10\% -- validation
\end{itemize}

Impartirea s-a facut la nivel de conversatii, cu seed fix.


\section{Hiperparametri utilizati}

\begin{itemize}
    \item modele: LLaMA3.1-8B-Instruct, RoLLaMA3.1-8B-Instruct
    \item LoRA: r=8, alpha=16, dropout=0.05
    \item quantizare: 4-bit (nf4)
    \item epoci: 2
    \item batch size efectiv: 8
    \item learning rate: 5e-5
    \item scheduler: cosine
    \item warmup ratio: 0.1
    \item max sequence length: 2048
    \item optimizer: paged AdamW 8-bit
    \item gradient checkpointing: activat
\end{itemize}


\section{Rezultate obtinute}

\subsection{Training loss LLaMA}
\begin{itemize}
    \item Step 10: 3.5843
    \item Step 20: 0.8830
    \item Step 30: 0.6452
    \item Step 40: 0.5084
    \item Step 50: 0.4188
\end{itemize}

\subsection{Training loss RoLLaMA}
\begin{itemize}
    \item Step 10: 2.2879
    \item Step 20: 0.8347
    \item Step 30: 0.6144
    \item Step 40: 0.4872
    \item Step 50: 0.3980
\end{itemize}


\subsection{Comportamentul modelelor}

\textbf{RoLLaMA}
\begin{itemize}
    \item respecta mai strict regulile din prompt
    \item pastreaza mai bine contextul conversatiei
    \item ofera raspunsuri mai stabile
    \item uneori limbajul este mai simplu si mai putin fluent fata de LLaMA
\end{itemize}

\textbf{LLaMA}
\begin{itemize}
    \item are un limbaj mult mai fluent si natural
    \item suna mai ``frumos'' si mai apropiat de o conversatie reala
    \item halucineaza mai usor
    \item genereaza mai des raspunsuri incoerente sau amuzante (``perle'')
\end{itemize}


\subsection{Exemple corecte}

\begin{itemize}
    \item \textbf{Student:} ``Ce problema aveti?'' \\
    \textbf{Pacient:} ``Ma doare gingia aici, in partea dreapta, mai ales cand ating zona.''
    
    \item \textbf{Student:} ``Aveti ceva umflatura?'' \\
    \textbf{Pacient:} ``Da, am o umflatura mare peste maseaua de minte, pe partea stanga.''
\end{itemize}


\subsection{Exemple gresite (asa zisele perle)}

\textbf{LLaMA:}
\begin{itemize}
    \item ``Am un dinte care pare mai inalt decat ceilalti si ma doare putin cand ma uit la el.''
    \item ``Am durere gingie si nu pot sa ma asez in scaun fara sa doara.''
\end{itemize}


\textbf{RoLLaMA:}
\begin{itemize}
    \item \textbf{Student:} ``Ai simtit ceva lichid cu gust ciudat?'' \\
    \textbf{Pacient:} ``Da, am gustat ceva lichid cu gust sarat, ca apa de gura, cand m-am spalat pe dinti.''
\end{itemize}


\end{document}
